\section{Model Description}\label{model-description}

A framework recently proposed outlines ten plausible futures for Earth-originating civilization 1000 years from now, structured around variations in governance type, resource regime, and systemic resilience \citep{Haqq-Misra2025-rt}. These scenarios are grounded in Earth's present-day sociotechnical conditions and extrapolated using techniques from futures studies, including PEST analysis (Political, Economic, Social, Technological), cross-consistency assessment, and structured scenario worldbuilding. Although the framework explores plausible civilizational trajectories, it does not explicitly model the temporal dynamics of collapse and recovery, nor their implications for the persistence and detectability of technosignatures. This leaves open a crucial question: how do different sociotechnical paths shape the \emph{ON}/\emph{OFF} dynamics of technosignature production (and potential detection)?

To explore this, we build a dynamic model that simulates long-term evolution of technological capacity and resource stocks under varying governance structures, resource regimes, and hazard exposure. Collapse and recovery cycles emerge naturally from feedback between growth and depletion, combined with stochastic existential risk. We use this model to explore how long civilizations remain detectable, how often they go silent, and what kinds of futures might dominate the galactic landscape. From our simulations, we compute the average duty cycle of each scenario, and from this, we calculate an effective detectability duration ($L_{\mathrm{eff}}$), allowing us to assess not just how long civilizations last, but how often we might catch them in the act of being detectable.

\subsection{Simulation framework}\label{simulation-framework}

We developed a hybrid deterministic-stochastic model to simulate the long-term behavior of technological civilizations, with the goal of capturing collapse-recovery dynamics and their effect on technosignature detectability. This model extends the scenario-based framework introduced by Haqq-Misra et al. \citep{Haqq-Misra2025-rt}, adding temporal dynamics and explicit feedback between resource depletion, technological growth, and systemic collapse. The simulation spans a fixed period of 1000 years, discretized into yearly time steps. At each step, a civilization is characterized by its technological capacity, \(T(t)\), and its available resource stock, \(R(t)\). Technological capacity increases linearly each year according to a fixed growth rate \(r\), which we take directly from the values reported in Haqq-Misra et al., Table 9. Simultaneously, the available resource stock is depleted at a constant rate \(\delta\). If the resource level falls to zero (i.e., \(R(t) < 0\)), or if an existential risk event occurs (represented as a probabilistic hazard drawn at each time step) the civilization undergoes a collapse. Collapse represents a systemic breakdown in infrastructure, coordination, and technological continuity. It leads to a sharp reduction in technological capacity, determined by a collapse factor \(c_{f}\), which defines the proportion of technology that survives the event. This loss can be interpreted as the result of disruptions to power grids, communication networks, knowledge retention, or the abandonment of complex systems.

Following collapse, the system enters a dormant recovery period of duration \(r_{d}\). During this phase, the society is unable to grow or consume resources, reflecting a temporary halt in technological and economic activity. This delay allows time for institutional reorganization, population stabilization, or environmental recovery. Once the recovery period concludes, the resource stock is partially replenished to a fraction \(r_{f}\) of its original value \(R_{0}\). This recovery fraction represents retained infrastructure, stored reserves, or ecological regeneration. A full recovery (\(r_{f} \approx 1\)) implies robust contingency planning and strong institutional memory, while lower values indicate more severe losses and diminished rebuilding capacity. This cycle of growth, collapse, dormancy, and recovery can repeat multiple times during the simulation window, depending on the fragility of the system, the frequency of stochastic hazards, and the depletion rate of resources.

Formally, let the initial conditions be \(T_{0} = 1\), and \(R_{0} \in {\mathbb{R}}_{+}\ \), and define the indicator function as \(\chi(t) \in \{ 0,1\}\), with \(\chi(t) = 1\) if the system is active (not collapsed or recovering) at time \(t\), and \(\chi(t) = 0\) otherwise. The update rules for each time step are given by
\begin{equation}
\begin{aligned}
T_{t+1} &= \chi_t \left( T_t + r \right) + \left(1 - \chi_t\right) T_t ,\\
R_{t+1} &= \chi_t \left( R_t - \delta \right) + \left(1 - \chi_t\right) R_t .
\end{aligned}
\label{eq:update}
\end{equation}

Equations~\eqref{eq:update}--\eqref{eq:recovery} define the full dynamical cycle of growth, collapse, and recovery in the model.

We include both a deterministic and a stochastic trigger for collapse to reflect two distinct failure modes: internal depletion and external shocks. Collapse is deterministically triggered when the resource stock is fully depleted (i.e.,\(\ R(t) \leq 0\)). It may also be triggered by a stochastic existential hazard, which we model as a Bernoulli process. At each time step \emph{t}, a uniform deviate \(U(t)\ \sim\ \ U(0,\ 1)\) is drawn. If \(U(t)\  < h\), where \(h \in \lbrack 0,1\rbrack\) is the hazard rate specific to the scenario, then a collapse is triggered independently of resource levels. 

Upon collapse, the system state is reset according to
\begin{equation}
\begin{aligned}
T_{t+1} &= c_f \, T_t ,\\
R_{t+1} &= 0 .
\end{aligned}
\label{eq:collapse}
\end{equation}

The system then remains inactive for $r_d$ time steps, after which resource availability is restored as
\begin{equation}
R_{t+r_d} = r_f \, R_0 .
\label{eq:recovery}
\end{equation}

We define a civilization as being ON when it maintains a positive level of resources (i.e., \(\ R(t) > 0\)). This assumption reflects the idea that some minimal material throughput is necessary for societal function and detectability, including energy consumption, communication, and technosignature production. Conversely, when resources are fully depleted (\(R(t) = 0\)), the civilization is considered OFF, representing a dormant, collapsed, or otherwise undetectable state. This binary definition underlies the calculation of duty cycle and related metrics discussed in subsection ~\ref{metrics}.

\subsection{Scenario typologies and parameter mapping}\label{scenario-typologies-and-parameter-mapping}

We apply our model to the ten future scenarios to simulate the long-term behavior of civilizations under diverse sociotechnical conditions. Each scenario describes a distinct vision for humanity's trajectory, defined by combinations of resource abundance, governance structure, institutional resilience, and technological integration. Our goal is to translate these qualitative narratives into a coherent set of quantitative parameters that govern dynamics of growth, collapse, and recovery. Scenarios are characterized along three conceptual axes: resource regime (scarcity vs. non-scarcity), governance type (from autocratic to participatory), and structural design (hierarchical vs. distributed). Scenarios are also mapped onto six thematic clusters (Table~\ref{tab:scenarios}).

\begin{table*}[htbp]
\centering
\caption{Scenario descriptions and sociotechnical tags}
\label{tab:scenarios}
\setlength{\tabcolsep}{8pt}
\renewcommand{\arraystretch}{1.25}
{\raggedright
\begin{tabular}{l p{4.5cm} p{5.5cm}}
\toprule
Scenario & Tags & Description \\
\midrule
S1: Big Brother is Watching &
Scarcity + Rule by One + Hierarchical \newline Cluster 6 (Earth 2024) &
Centralized authoritarian system under resource stress \\

S2: Wild West &
Scarcity + Rule by Few + Hierarchical \newline Cluster 6 (Earth 2024) &
Deregulated, competitive, and fragmented society \\

S3: Golden Age &
Non-scarcity + Rule by All + Distributed \newline Cluster 6 (Earth 2024)$^{\dagger}$ &
Stable egalitarian society with abundant resources \\

S4: Living with the Land &
Scarcity + Rule by Few + Hierarchical \newline Cluster 5 (Simplicity) &
Low-impact, cyclical, ecologically constrained society \\

S5: Transhumanism &
Non-scarcity + Rule by All + Distributed \newline Cluster 4 (Engineering) &
Technologically integrated, post-biological civilization \\

S6: Sword of Damocles &
Scarcity + Rule by Few + Hierarchical \newline Cluster 4 (Engineering) &
High-tech and high-risk system with cascading failure potential \\

S7: Restoration &
Non-scarcity$^{\ddagger}$ + Rule by All + Distributed \newline Cluster 3 (Sustainability) &
Recovered post-collapse world with resilience \\

S8: Ouroboros &
Scarcity + Rule by Few + Hierarchical \newline Cluster 3 (Sustainability) &
Oscillatory civilization with cyclical collapse and regrowth \\

S9: Deus Ex Machina &
Non-scarcity + Rule by All + Distributed \newline Cluster 2 (Tech Escape) &
Aggressive technological expansion with instability risk \\

S10: Out of Eden &
Non-scarcity + Rule by All + Distributed \newline Cluster 1 (Parallel Life) &
Harmonized coexistence with nature and machines \\
\bottomrule
\end{tabular}
}

\smallskip
{\footnotesize $^{\dagger}$Cluster assignments follow \citet{Haqq-Misra2025-rt}; S3 shares the Earth~2024 technology cluster with S1 and S2 despite differing in global factors.\\
$^{\ddagger}$S7 is tagged non-scarcity in the original framework but parameterized with scarcity-level resource pressure to model post-collapse recovery (see text).}
\end{table*}

% \begin{sidewaystable*}[p]
% \centering
% \caption{Scenario descriptions and sociotechnical tags}
% \label{tab:scenarios}
% \setlength{\tabcolsep}{8pt}
% \renewcommand{\arraystretch}{1.25}
% {\raggedright
% \begin{tabular}{l l p{0.3\textwidth} p{0.4\textwidth}}
% \toprule
% Scenario & Name & Tags & Description \\
% \midrule
% S1 & Big Brother is Watching &
% Scarcity + Rule by One + Hierarchical \newline Cluster 6 (Earth 2024) &
% Centralized authoritarian system under resource stress \\

% S2 & Wild West &
% Scarcity + Rule by Few + Hierarchical \newline Cluster 6 (Earth 2024) &
% Deregulated, competitive, and fragmented society \\

% S3 & Golden Age &
% Non-scarcity + Rule by All + Distributed \newline Cluster 6 (Earth 2024) &
% Stable egalitarian society with abundant resources \\

% S4 & Living with the Land &
% Scarcity + Rule by Few + Hierarchical \newline Cluster 5 (Simplicity) &
% Low-impact, cyclical, ecologically constrained society \\

% S5 & Transhumanism &
% Non-scarcity + Rule by All + Distributed \newline Cluster 4 (Engineering) &
% Technologically integrated, post-biological civilization \\

% S6 & Sword of Damocles &
% Scarcity + Rule by Few + Hierarchical \newline Cluster 4 (Engineering) &
% High-tech and high-risk system with cascading failure potential \\

% S7 & Restoration &
% Non-scarcity + Rule by All + Distributed \newline Cluster 3 (Sustainability) &
% Recovered post-collapse world with resilience \\

% S8 & Ouroboros &
% Scarcity + Rule by Few + Hierarchical \newline Cluster 3 (Sustainability) &
% Oscillatory civilization with cyclical collapse and regrowth \\

% S9 & Deus Ex Machina &
% Non-scarcity + Rule by All + Distributed \newline Cluster 2 (Tech Escape) &
% Aggressive technological expansion with instability risk \\

% S10 & Out of Eden &
% Non-scarcity + Rule by All + Distributed \newline Cluster 1 (Parallel Life) &
% Harmonized coexistence with nature and machines \\
% \bottomrule
% \end{tabular}
% }
% \end{sidewaystable*}

We assign parameter values based on the sociotechnical features of each scenario. The growth rate \(r\) is taken directly from Haqq-Misra et al. (2025) and held fixed. We note that assuming technological capacity increases linearly each year via \(r\) is a simplification, especially given scenarios that imply discontinuities such as collapse points or radical transformation points; we retain it here for consistency with the source parameterization. The remaining parameters (initial resource stock \(R_{0}\), per-year depletion rate \(\delta\), collapse depth \(c_{f}\), recovery delay \(r_{d}\), recovery fraction \(r_{f}\), and existential hazard rate \(h\)) are derived from the narrative structure of each scenario and converted into quantitative values through structured reasoning. 

We assume that governance structure influences how technological civilizations respond to systemic stress and collapse. Highly centralized systems (rule by one) could be more vulnerable to severe and prolonged disruptions, as decision-making and critical infrastructure are concentrated and potentially less adaptable. When collapse occurs (due to resource exhaustion or external hazard) such civilizations may face steeper declines in technological capacity and slower recovery, particularly if institutional continuity is fragile. In contrast, more distributed governance (rule by all) might support greater redundancy and flexibility, enabling shallower collapses and faster recovery following disruption. To reflect these hypothesized patterns, we selected representative parameter value ranges for collapse depth \(c_{f}\), recovery delay \(r_{d}\), and recovery fraction \(r_{f}\) that align with the governance structure assumed in each scenario (Table~\ref{tab:governance}). These values are not based on empirical observation, but represent internally consistent assumptions designed to explore contrasting trajectories in our scenario-based simulations. 

One exception is the Restoration scenario (S7). While it is categorized as non-scarcity and governed by distributed institutions, it is conceptually framed as a rebuilding civilization emerging from an earlier collapse. To reflect this, we initialize the simulation with a reduced effective resource stock and a moderate depletion rate (parameters that would otherwise suggest scarcity), thereby modeling a society recovering from prior systemic failure. This allows us to represent S7 as a post-collapse recovery regime without forcing a collapse to occur within the simulation window.

The Ouroboros scenario (S8) is treated differently. In the original framework, S8 is described as a civilization that repeatedly undergoes expansion, collapse, and regrowth, with the year 3000 situated partway through an ongoing recovery phase. Rather than imposing externally timed or periodic collapse events, we model this behavior as an emergent dynamical regime. Specifically, S8 is assigned moderate technological growth, partial technological survival during collapse, and a relatively high recovery fraction of planetary resources. This combination allows repeated collapse--recovery cycles to arise endogenously from the interaction between resource depletion and rebuilding, and permits exploration of whether oscillatory behavior remains stable or dampens over long timescales.

Similarly, Deus Ex Machina (S9) is tagged as non-scarcity in the original framework but is parameterized with moderate resource pressure ($R_0 = 900$, $\delta = 1.3$) to reflect the high extraction demands of its aggressive technological expansion program, which drives resource depletion despite the scenario's nominally post-scarcity classification.

\begin{table}[htbp]
\centering
\caption{Mapping sociotechnical structure to parameter logic}
\label{tab:governance}
\begin{tabular}{llll}
\toprule
Governance Type & Collapse Depth $c_f$ & Recovery Delay $r_d$ & Recovery Fraction $r_f$ \\
\midrule
Rule by One & 0.5--0.7 & 40--60 & 0.10--0.20 \\
Rule by Few & 0.2--0.6 & 50--100 & 0.10--0.35 \\
Rule by All$^{*}$ & 0.3--1.0 & 0--70 & 0.20--1.0 \\
\bottomrule
\end{tabular}

\smallskip
{\footnotesize $^{*}$Post-scarcity idealizations (S3, S10) use $c_f = 1.0$, $r_d = 0$, $r_f = 1.0$, reflecting no-collapse conditions. S7 and S9 use parameters outside the typical governance range as narrative exceptions (see text).}
\end{table}

Scenarios also differ fundamentally in their assumptions about planetary resource conditions. Rather than define scarcity solely in terms of starting resources, we interpret it as a function of resource pressure (i.e., how quickly a civilization depletes its reserves). That is, systems with a large initial stock may still experience scarcity if it depletes resources at a high rate (e.g., S6). Conversely, a civilization with a modest stock may persist through repeated collapse--recovery cycles if post-collapse resource recovery partially replenishes its base (e.g., S8, S9), or may rebuild following an earlier systemic failure (e.g., S7).

\begin{table}[htbp]
\centering
\caption{Interpreting resource pressure from $R_0$ and $\delta$}
\label{tab:resources}
\begin{tabular}{llll}
\toprule
Resource Pressure & $R_0$ & $\delta$ & Scenarios \\
\midrule
High & $\leq 1000$ & 1.0--2.5 & S1, S2, S4, S7, S8, S9 \\
Moderate & 400--1600 & 0.5--1.0 & S6 \\
Low & $\geq 1200$ & 0.0--0.5 & S3, S5, S10 \\
\bottomrule
\end{tabular}
\end{table}

We implement Monte Carlo simulations to explore how uncertainty in model parameters shapes the long-term behavior of each scenario. Each scenario's baseline values (Table~\ref{tab:params}) are perturbed using probabilistic variation to capture plausible uncertainty (Table~\ref{tab:uncertainty}). Initial resource stock (\(R_{0}\)) and depletion rate (\(\delta\)) are sampled from normal distributions with modest standard deviations to represent uncertainty in ecological baselines and consumption patterns. Collapse depth (\(c_{f}\)) and recovery fraction (\(r_{f}\)) are sampled using triangular distributions centered on scenario-specific values, allowing outcomes to concentrate near expected levels while still capturing asymmetries in resilience and recovery capacity. Recovery delay (\(r_{d}\)) is sampled from a normal distribution with a standard deviation of $\pm 10$ years, reflecting variability in rebuilding pace across different societal structures.

The existential hazard rate (\(h\)) is held fixed for each scenario because it encodes an assumed background level of external risk (such as AI misalignment, interstellar conflict, or catastrophic environmental feedbacks) rather than internal system dynamics. Unlike other parameters whose variability arises from structural or operational uncertainty, \(h\) reflects a narrative-level judgment about a civilization's exposure to exogenous threats. Scenarios characterized by unstable geopolitical contexts, risky technological dependencies, or unresolved existential challenges (e.g., S1, S6, S7) are assigned higher hazard rates to reflect this vulnerability, while S5 carries a small residual risk due to high systemic complexity. In contrast, scenarios with societal stability, low-impact lifestyles, or resilient infrastructure (e.g., S3, S4, S8, S9, S10) are assigned zero or near-zero hazard rates. Keeping \(h\) fixed rather than sampling it probabilistically allows for clearer attribution of collapse outcomes to either endogenous fragility or exogenous shocks, rather than conflating the two.

\begin{table*}[htbp]
\centering
\caption{Final parameter values for each scenario}
\label{tab:params}
\begin{tabular}{lccccccc}
\toprule
Scenario & $r$ & $R_0$ & $\delta$ & $c_f$ & $r_d$ & $r_f$ & $h$ \\
\midrule
S1 & 0.0020 & 400 & 1.2 & 0.6 & 40 & 0.15 & 0.003 \\
S2 & 0.0026 & 700 & 1.5 & 0.3 & 70 & 0.10 & 0.001 \\
S3 & 0.0001 & 2400 & 0.0 & 1.0 & 0 & 1.0 & 0.0 \\
S4 & 0.0060 & 600 & 2.5 & 0.2 & 100 & 0.10 & 0.0 \\
S5 & 0.0003 & 2500 & 0.1 & 0.9 & 20 & 0.80 & 0.0005 \\
S6 & 0.0018 & 1600 & 1.0 & 0.4 & 50 & 0.25 & 0.005 \\
S7 & 0.0060 & 400 & 1.5 & 0.3 & 60 & 0.20 & 0.001 \\
S8 & 0.0015 & 800 & 1.3 & 0.6 & 70 & 0.35 & 0.0 \\
S9 & 0.0011 & 900 & 1.3 & 0.5 & 70 & 0.35 & 0.0 \\
S10 & 0.0002 & 2700 & 0.0 & 1.0 & 0 & 1.0 & 0.0 \\
\bottomrule
\end{tabular}
\end{table*}

\begin{table}[htbp]
\centering
\caption{Parameter uncertainty assumptions for Monte Carlo sampling}
\label{tab:uncertainty}
\begin{tabular}{lll}
\toprule
Parameter & Distribution & Uncertainty \\
\midrule
$R_0$ & Normal & $\pm 5\%$ of mean \\
$\delta$ & Normal & $\pm 5\%$ of mean \\
$c_f$ & Triangular & $\pm 5\%$ of mode \\
$r_d$ & Normal & $\pm 10$ years \\
$r_f$ & Triangular & $\pm 5\%$ of mode \\
\bottomrule
\end{tabular}
\end{table}

\FloatBarrier

\subsection{Metrics}\label{metrics}

We define four metrics to reflect different aspects of technospheric persistence and collapse dynamics across Monte Carlo simulations. These metrics are designed to capture the onset, frequency, and continuity of civilization-level activity, with a focus on their implications for detectability.

The first metric, \(T_{c}\), is the mean time to first collapse across replicates. It serves as a proxy for system fragility: scenarios with low values of \(T_{c}\) tend to break down early, while those with higher values demonstrate greater initial stability. When a simulation run does not experience collapse within the 1000-year window, the full timespan is used. The second metric, \(f_{nc}\), denotes the fraction of runs that avoid collapse entirely. This value reflects the overall robustness of a scenario and helps differentiate consistently stable futures from those vulnerable to systemic failure. We also measure \(D_{c}\), the average duty cycle, defined as the fraction of time a civilization remains active, that is, maintaining nonzero resource levels and thus potentially emitting technosignatures. This definition assumes that minimal resource availability is a prerequisite for detectable activity. Scenarios with high \(D_{c}\) maintain persistent activity, while those with low \(D_{c}\) may cycle through long dormant periods or experience early termination. Finally, \(N_{c}\), the mean number of collapse-recovery cycles, captures whether a scenario tends to exhibit singular collapse events or undergo repeated disruptions over time. Together, these metrics provide a compact but informative summary of each scenario's trajectory, enabling comparisons across sociotechnical regimes and offering insight into the reliability and intermittency of technosignature emission.
